\section{Exercici 2}
Siguin $z_1, z_2$ les arrels del polinomi
\[
x^2 + ax + b = 0, \quad a,b \in \mathbb{R}.
\]

Com que els coeficients són reals, si les arrels no són reals, són conjugades:
\[
z_2 = \overline{z_1}.
\]

Per a qualsevol enter positiu $n$:
\[
z_1^n + z_2^n = z_1^n + \overline{z_1}^n.
\]

Com que
\[
\overline{z_1^n} = \overline{z_1}^n,
\]
tenim:
\[
z_1^n + z_2^n = z_1^n + \overline{z_1^n} = 2 \,\mathrm{Re}(z_1^n).
\]

Per tant,
\[
z_1^n + z_2^n \in \mathbb{R}.
\]

Resolem el cas concret $x^2 - 2x + 2 = 0$:
\[
x = \frac{2 \pm \sqrt{4 - 8}}{2} = \frac{2 \pm \sqrt{-4}}{2} = 1 \pm i.
\]

Així,
\[
z_1 = 1+i, \quad z_2 = 1-i.
\]

Forma polar:
\[
|z_1| = \sqrt{1^2 + 1^2} = \sqrt{2}, \quad \arg(z_1) = \frac{\pi}{4}.
\]

\[
z_1 = \sqrt{2}\left(\cos \frac{\pi}{4} + i \sin \frac{\pi}{4}\right),
\]
\[
z_2 = \sqrt{2}\left(\cos \frac{\pi}{4} - i \sin \frac{\pi}{4}\right).
\]

Potències:
\[
z_1^n = (\sqrt{2})^n \left(\cos \frac{n\pi}{4} + i \sin \frac{n\pi}{4}\right),
\]
\[
z_2^n = (\sqrt{2})^n \left(\cos \frac{n\pi}{4} - i \sin \frac{n\pi}{4}\right).
\]

Sumant:
\[
z_1^n + z_2^n = 2(\sqrt{2})^n \cos\left(\frac{n\pi}{4}\right).
\]

Resultat final:
\[
\boxed{
z_1^n + z_2^n = 2(\sqrt{2})^n \cos\left(\frac{n\pi}{4}\right)
}
\]